% !TEX root = userguide.tex

\chapter{External routines}
\pgst{tfemch}

\label{ch:external}

\def\chaptertitle{External routines}
\def\headdate{11th March 2016}

In this chapter the \texttt{external} tree in \tfem\ is briefly described. The \texttt{external} tree
does \emph{not} contain the routines, only the Makefiles for compiling \& building
libraries. The routines itself require licenses and need to be obtained
separately.

\section{HSL}
\label{sec:hsl_external}

The default solvers in \tfem\ are MA41 and MA57 from the HSL library (see Sec.~\ref{sec:hslsparse}).
Within the department most students use the routines from HSL2002, since these have been licensed. However, MA57 does not have scaling and metis renumbering (see Sec.~\ref{sec:metis5}) in that version. A more recent version is needed. 

HSL packages have a free personal license for academics. 
See the file \texttt{external/README} and \texttt{external/hsl/README} for more information. 

Provided the routines have been obtained, the recent MA41 and MA57 routines can be compiled and added to the solver library by
typing
\begin{quote}
   \texttt{make libhsl2}
\end{quote}
from the \tfem\ root directory. See the root \texttt{Makefile} which HSL packages are required.

The add-on \texttt{hsl\_extra} (see Sec.~\ref{sec:hsl_extra_addon}) requires additional routines from the HSL library. Provided the routines have been obtained, the routines can be compiled and added to the solver library by
typing
\begin{quote}
   \texttt{make libhsl3}
\end{quote}
from the \tfem\ root directory. See the root \texttt{Makefile} which HSL packages are required. Note, that \texttt{libhsl3} requires \texttt{libhsl2}.

\section{Numerical Recipes}
\label{sec:nr}

The add-on particle\_tracking (see Sec.~\ref{sec:particletrackingaddon}) needs some (modified for double precision) routines from the Numerical Recipes library\cite{Press92,Press96}, available
at \href{http://www.nr.com}{www.nr.com}.
The software is not free and need to be purchased separately or as part of a book.
See the file \texttt{external/README} and \texttt{external/nr/README} for more information. Provided the routines have been obtained (and modified for double precision) the routines\footnote{See the \texttt{Makefile} for which routines are needed.} can be compiled and added to the solver library by
typing
\begin{quote}
   \texttt{make nr}
\end{quote}
from the \tfem\ root directory.

\section{Sparskit2}
\label{sec:sparskit2}

The add-on \texttt{sk\_solve} (see Sec.~\ref{sec:sksolve}) needs some routines from the Sparskit2 
library\footnote{\href{http://www-users.cs.umn.edu/~saad/software/SPARSKIT}{http://www-users.cs.umn.edu/\textasciitilde{}saad/software/SPARSKIT}}. 
The software is free (LGPL). See the file \texttt{external/README} and \texttt{external/sparskit2/README} for more information. Provided the routines have been obtained, the routines can be compiled and added to the solver library by
typing
\begin{quote}
   \texttt{make sparskit2}
\end{quote}
from the \tfem\ root directory.

\section{Removing external packages}
\label{sec:external_removing}

It is recommended not to remove/clean individual external packages.
Clean the complete external tree at once using
\begin{quote}
   \texttt{make externalclean}
\end{quote}
instead of cleaning individual packages. The reason is that cleaning an
    individual package will also remove the complete solver library, leaving
    other compiled packages ``hanging'', unable to build them again. 

Similarly, if you want to go back to \texttt{libhsl} from \texttt{libhsl2} do a
clean of the complete external tree:
\begin{quote}
      \texttt{make externalclean}
\end{quote}
    then build libhsl
\begin{quote}
      \texttt{make libhsl}
\end{quote}
    and rebuild non-HSL external packages again, if needed.






